% =====================================================================
%  MODELLO DI ESERCIZIO PER LA BANCA — copialo in
%     Banca_Esercizi/<materia>/<argomento>_NN.tex
%  e aggiungilo a Banca_Esercizi/catalogo_esercizi.tex.
%  Un file = un esercizio. Si usa sia nelle verifiche sia negli
%  esercizi a casa (lì vale solo la fila A e i punti non si vedono).
% ---------------------------------------------------------------------
%  Argomento:            Livello: base / intermedio / avanzato
%  Tempo stimato: .. min
% =====================================================================
\begin{esercizio}          % [4] dopo \begin{esercizio} = punteggio unico
% --- DATI: nome, valore fila A, fila B, fila C (unità base: 2.2e3 = 2,2 kΩ)
%     Scegli valori che NON diano risultati tondi: nelle soluzioni i
%     risultati con 2 cifre significative o meno vengono segnati con [!].
\dato{V}{12}{15}{18}
\dato{R}{390}{560}{470}
% --- CALCOLI: si fanno da soli per le tre file
%     [nome da mostrare]{nome}{formula con \x{dato}}
\calcola[I]{I}{\x{V}/\x{R}}

% --- TESTO
Un resistore da \mis{R}{\ohm} è alimentato con \mis{V}{\volt}.
\begin{parts}
  \part[2] Calcola la corrente.
\end{parts}

% --- CIRCUITI: generatore con simbolo europeo e segni ai morsetti
%   \draw (0,0) to[vsource, l=$V$, n=G] (0,3) ... ;   \morsetti{G}
%   (il + va verso il punto di arrivo del "to")

% --- SOLUZIONE (soluzioni per studenti, griglia di correzione, catalogo)
%     \criterio{...} compare solo nella griglia di correzione, mai agli studenti
\begin{solution}
  $I = V/R = \val{V}/\val{R} \approx \mis{I}{\ampere}$
  \criterio{1 pt formula, 1 pt risultato con unità.}
\end{solution}
\end{esercizio}
